\documentclass[11pt]{article}

\usepackage[margin=1in]{geometry}
\usepackage{amsmath,amssymb,amsthm,mathtools}
\usepackage{lmodern}

\title{Random Intersections of Planes Generated by Uniform Points on the Sphere}
\author{}
\date{}

\begin{document}
\maketitle

\section*{Problem}

Let $U$ be uniform on the unit sphere $S^2 \subset \mathbb{R}^3$. Take three iid copies
$x_1,x_2,x_3 \sim U$. Almost surely these determine a unique plane, which we write in oriented
form as
\[
P = \{x \in \mathbb{R}^3 : n \cdot x = h\},
\qquad n \in S^2,\quad h \in [-1,1].
\]

Now take three independent such planes
\[
P_i = \{x : n_i \cdot x = h_i\}, \qquad i=1,2,3,
\]
constructed from disjoint triples of iid uniform points on $S^2$. Almost surely the planes meet at
a unique point
\[
X = P_1 \cap P_2 \cap P_3.
\]
We want the distribution of
\[
R = \|X\|.
\]

\subsection*{Geometric Intuition}

The problem combines two geometric phenomena. First, three random points on $S^2$ typically lie on
a circle (the intersection of $S^2$ with a plane), and the density of planes is not uniform: planes
close to the equator (where $|h|$ is small) are much more common than planes near the poles. This
explains the factor $(1-h^2)$ in the offset density $g(h)$.

Second, three planes in general position meet at a single point. The randomness in the three normals
$n_1, n_2, n_3$ introduces randomness in both the location and distance of this intersection from
the origin. The isotropy of the problem (invariance under rotations of $\mathbb{R}^3$) means that
the distance $R$ depends only on the norms, not on the directions, which is why the problem reduces
to a radial density $f_R(r)$.

\section*{Summary of the Result}

The key exact fact is the law of one random plane:

\begin{itemize}
\item its oriented normal $n$ is uniform on $S^2$,
\item $n$ is independent of its signed offset $h$,
\item the offset density is
\[
g(h) = \frac{3}{4}(1-h^2), \qquad -1 \le h \le 1.
\]
\end{itemize}

From this, the intersection point density is
\[
f_X(x)
=
\mathbb{E}\!\left[
\left|\det(n_1,n_2,n_3)\right|
\prod_{i=1}^3 g(n_i \cdot x)
\right].
\]
Hence
\[
f_X(x)
=
\left(\frac{3}{4}\right)^3
\mathbb{E}\!\left[
\left|\det(n_1,n_2,n_3)\right|
\prod_{i=1}^3 \bigl(1-(n_i \cdot x)^2\bigr)_+
\right],
\]
where $(a)_+ := \max(a,0)$.

By isotropy, $f_X(x)$ depends only on $\|x\|$, and therefore
\[
f_R(r) = 4\pi r^2 f_X(re_3), \qquad r \ge 0.
\]

I do not presently have a closed elementary form for $f_R$, but the formulas below reduce it to an
explicit low-dimensional integral.

\section*{1. Law of a Single Random Plane}

Let $\sigma$ denote surface area measure on $S^2$, and let $\mu = \sigma/(4\pi)$ be the uniform
probability measure.

Take three iid uniform points $x_1,x_2,x_3 \in S^2$. Almost surely they determine a unique plane,
which can be written as
\[
P = \{x : n \cdot x = h\}.
\]

For fixed $n \in S^2$ and $h \in (-1,1)$, define
\[
\rho = \sqrt{1-h^2}.
\]
Then the intersection
\[
C(n,h) := S^2 \cap \{x : n \cdot x = h\}
\]
is a circle of radius $\rho$.

Choose a rotation $R_n \in SO(3)$ with $R_n e_3 = n$, and parametrize the circle by
\[
x_i = h\,n + \rho\,R_n v(\theta_i),
\qquad
v(\theta) = (\cos\theta,\sin\theta,0),
\qquad i=1,2,3.
\]
Thus
\[
(n,h,\theta_1,\theta_2,\theta_3) \mapsto (x_1,x_2,x_3)
\]
parametrizes almost every non-collinear triple on $S^2$.

By rotational invariance it suffices to compute the Jacobian at $n=e_3$. Put
\[
v_i = v(\theta_i), \qquad w_i = (-\sin\theta_i,\cos\theta_i,0).
\]

Use local coordinates on $S^2$ near $e_3$ given by infinitesimal rotations about the $e_1$- and
$e_2$-axes, with parameters $a,b$. Then the six tangent vectors in $(\mathbb{R}^3)^3$ are
\[
\partial_{\theta_i}x
=
(0,\dots,0,\rho w_i,0,\dots,0),
\]
\[
\partial_h x
=
\bigl(-\tfrac{h}{\rho}v_1+e_3,\,-\tfrac{h}{\rho}v_2+e_3,\,-\tfrac{h}{\rho}v_3+e_3\bigr),
\]
\[
\partial_a x
=
(e_1 \times x_1,\ e_1 \times x_2,\ e_1 \times x_3),
\]
\[
\partial_b x
=
(e_2 \times x_1,\ e_2 \times x_2,\ e_2 \times x_3).
\]

Let $G$ be the $6 \times 6$ Gram matrix of these vectors. A direct Schur-complement computation
gives
\[
\det G = \rho^6 \det S,
\]
where
\[
S =
\begin{pmatrix}
3/\rho^2 & (\sum \sin\theta_i)/\rho & -(\sum \cos\theta_i)/\rho\\
(\sum \sin\theta_i)/\rho & \sum \sin^2\theta_i & -\sum \sin\theta_i\cos\theta_i\\
-(\sum \cos\theta_i)/\rho & -\sum \sin\theta_i\cos\theta_i & \sum \cos^2\theta_i
\end{pmatrix}.
\]
This can be factorized as
\[
S
=
\operatorname{diag}(\rho^{-1},1,1)\,
ZZ^\top\,
\operatorname{diag}(\rho^{-1},1,1),
\]
with
\[
Z =
\begin{pmatrix}
1 & 1 & 1\\
\sin\theta_1 & \sin\theta_2 & \sin\theta_3\\
-\cos\theta_1 & -\cos\theta_2 & -\cos\theta_3
\end{pmatrix}.
\]
Hence
\[
\det G = \rho^4 (\det Z)^2.
\]

Now
\[
\frac{1}{2}|\det Z|
\]
is the Euclidean area of the triangle in the plane with vertices
\[
(\cos\theta_i,\sin\theta_i), \qquad i=1,2,3.
\]
Call that area $A(\theta_1,\theta_2,\theta_3)$. Therefore the Jacobian is
\[
J = \sqrt{\det G} = 2\rho^2 A(\theta_1,\theta_2,\theta_3).
\]

Thus the surface-area decomposition is
\[
d\sigma(x_1)\,d\sigma(x_2)\,d\sigma(x_3)
=
2(1-h^2)\,A(\theta_1,\theta_2,\theta_3)\,
d\sigma(n)\,dh\,d\theta_1\,d\theta_2\,d\theta_3.
\]

After integrating out the circle angles, the dependence on $h$ is exactly $1-h^2$. Therefore the
law of $(n,h)$ has density
\[
\mathbb{P}\bigl((n,h)\in d\sigma\,dh\bigr)
=
c(1-h^2)\,d\sigma(n)\,dh
\]
for some constant $c$. Since
\[
\int_{S^2} d\sigma = 4\pi,
\qquad
\int_{-1}^1 (1-h^2)\,dh = \frac{4}{3},
\]
normalization gives
\[
c = \frac{3}{16\pi}.
\]

So the exact one-plane law is
\[
\boxed{
\mathbb{P}\bigl((n,h)\in d\sigma\,dh\bigr)
=
\frac{3}{16\pi}(1-h^2)\,d\sigma(n)\,dh
}
\]
or equivalently
\[
\boxed{
n \sim \text{Unif}(S^2),\qquad
n \perp h,\qquad
g(h) = \frac{3}{4}(1-h^2)\mathbf{1}_{[-1,1]}(h).
}
\]

\section*{2. Density of the Intersection Point}

Let
\[
N =
\begin{pmatrix}
n_1^\top\\
n_2^\top\\
n_3^\top
\end{pmatrix},
\qquad
H =
\begin{pmatrix}
h_1\\
h_2\\
h_3
\end{pmatrix}.
\]
Then almost surely
\[
X = N^{-1}H.
\]

Condition on $N$. Since $h_1,h_2,h_3$ are independent with density $g$,
\[
f_{X\mid N}(x\mid N)
=
|\det N|\prod_{i=1}^3 g(n_i\cdot x).
\]
Taking expectation in $N$ gives
\[
\boxed{
f_X(x)
=
\mathbb{E}\!\left[
\left|\det(n_1,n_2,n_3)\right|
\prod_{i=1}^3 g(n_i\cdot x)
\right].
}
\]

Substituting the explicit one-plane density,
\[
\boxed{
f_X(x)
=
\left(\frac{3}{4}\right)^3
\mathbb{E}\!\left[
\left|\det(n_1,n_2,n_3)\right|
\prod_{i=1}^3 \bigl(1-(n_i\cdot x)^2\bigr)_+
\right].
}
\]

This is already an exact density formula.

\section*{3. Radial Reduction}

By isotropy, $f_X(x)$ depends only on $r=\|x\|$, so we may set $x = re_3$.

Write
\[
t_i := n_i \cdot e_3 \in [-1,1],
\qquad
a_i := \sqrt{1-t_i^2},
\]
and parametrize
\[
n_i = (a_i\cos\phi_i, a_i\sin\phi_i, t_i).
\]
Since $d\sigma = dt\,d\phi$, one gets
\[
f_X(re_3)
=
\frac{27}{2048\pi^2}
\int_{[-1,1]^3}
\Bigl(\prod_{i=1}^3 (1-r^2 t_i^2)_+\Bigr)
K(t_1,t_2,t_3)\,dt_1dt_2dt_3,
\]
where
\[
K(t_1,t_2,t_3)
=
\int_0^{2\pi}\int_0^{2\pi}
|D(u,v;t_1,t_2,t_3)|\,du\,dv,
\]
with
\[
D(u,v;t_1,t_2,t_3)
=
a_1a_2t_3\sin u
-a_1a_3t_2\sin v
+a_2a_3t_1\sin(v-u).
\]

Therefore
\[
\boxed{
f_R(r) = 4\pi r^2 f_X(re_3), \qquad r\ge 0.
}
\]
That is,
\[
\boxed{
f_R(r)
=
\frac{27\,r^2}{512\pi}
\int_{[-1,1]^3}
\Bigl(\prod_{i=1}^3 (1-r^2 t_i^2)_+\Bigr)
K(t_1,t_2,t_3)\,dt_1dt_2dt_3.
}
\]

\section*{4. One More Angular Reduction}

For fixed $v$, the quantity $D(u,v)$ is of the form
\[
D(u,v) = A(v)\sin u + B(v)\cos u + C(v),
\]
where
\[
A(v)=a_1a_2t_3-a_2a_3t_1\cos v,
\]
\[
B(v)=a_2a_3t_1\sin v,
\]
\[
C(v)=-a_1a_3t_2\sin v.
\]
Let
\[
R(v) = \sqrt{A(v)^2+B(v)^2}.
\]
Then
\[
\frac{1}{2\pi}\int_0^{2\pi}|D(u,v)|\,du = \Psi(R(v),C(v)),
\]
where
\[
\Psi(R,C)
=
\begin{cases}
|C|, & |C|\ge R,\\[6pt]
\dfrac{2}{\pi}\left(\sqrt{R^2-C^2}+|C|\arcsin\!\dfrac{|C|}{R}\right), & |C|<R.
\end{cases}
\]
Hence
\[
K(t_1,t_2,t_3)
=
2\pi\int_0^{2\pi}\Psi(R(v),C(v))\,dv.
\]

So the radial density can be written as
\[
\boxed{
f_R(r)
=
\frac{27\,r^2}{256}
\int_{[-1,1]^3}
\Bigl(\prod_{i=1}^3 (1-r^2 t_i^2)_+\Bigr)
\left[
\int_0^{2\pi}\Psi(R(v),C(v))\,dv
\right]
dt_1dt_2dt_3.
}
\]

\section*{5. Two Rigorous Asymptotic Results}

The integral representation can be pushed further in two directions: the behavior near the origin,
and the order of the tail at infinity.

\subsection*{5.1. Behavior near the Origin}

At $x=0$ the exact density simplifies to
\[
f_X(0)
=
\left(\frac{3}{4}\right)^3 \mathbb{E}\!\left[\left|\det(n_1,n_2,n_3)\right|\right].
\]
So the first task is to compute the expectation of the random absolute determinant.

Fix $n_1,n_2$. Then
\[
\left|\det(n_1,n_2,n_3)\right|
=
\|n_1\times n_2\|
\left|
\frac{n_1\times n_2}{\|n_1\times n_2\|}\cdot n_3
\right|.
\]
Conditionally on $n_1,n_2$, the unit vector
\[
\frac{n_1\times n_2}{\|n_1\times n_2\|}
\]
is deterministic, while $n_3$ is uniform on $S^2$. Hence
\[
\mathbb{E}\!\left[
\left.
\left|
\frac{n_1\times n_2}{\|n_1\times n_2\|}\cdot n_3
\right|
\right| n_1,n_2
\right]
=
\frac{1}{2},
\]
because for a uniform point on $S^2$, the projection onto any fixed axis is uniform on $[-1,1]$.
Therefore
\[
\mathbb{E}\!\left[\left|\det(n_1,n_2,n_3)\right|\right]
=
\frac{1}{2}\,\mathbb{E}\!\left[\|n_1\times n_2\|\right].
\]

If $\Theta$ is the angle between $n_1$ and $n_2$, then $\cos\Theta$ is uniform on $[-1,1]$, and
\[
\|n_1\times n_2\| = \sin\Theta = \sqrt{1-\cos^2\Theta}.
\]
Hence
\[
\mathbb{E}\!\left[\|n_1\times n_2\|\right]
=
\frac{1}{2}\int_{-1}^1 \sqrt{1-t^2}\,dt
=
\frac{\pi}{4}.
\]
It follows that
\[
\boxed{
\mathbb{E}\!\left[\left|\det(n_1,n_2,n_3)\right|\right] = \frac{\pi}{8}.
}
\]

Consequently,
\[
\boxed{
f_X(0)=\left(\frac{3}{4}\right)^3 \frac{\pi}{8}=\frac{27\pi}{512}.
}
\]

Since the density formula is smooth in $x$ near the origin, we obtain
\[
f_X(x)=\frac{27\pi}{512}+O(\|x\|^2)
\qquad (\|x\|\to 0),
\]
and therefore
\[
\boxed{
f_R(r)=4\pi r^2 f_X(re_3)
=
\frac{27\pi^2}{128}r^2+O(r^4)
\qquad (r\to 0).
}
\]

\subsection*{5.2. Tail Order at Infinity}

Starting from
\[
f_R(r)
=
\frac{27\,r^2}{512\pi}
\int_{[-1,1]^3}
\Bigl(\prod_{i=1}^3 (1-r^2 t_i^2)_+\Bigr)
K(t_1,t_2,t_3)\,dt_1dt_2dt_3,
\]
we note that for large $r$, the factor $(1-r^2 t_i^2)_+$ restricts the integral to
$|t_i|\le r^{-1}$. Make the change of variables
\[
t_i=\frac{s_i}{r}.
\]
Then
\[
f_R(r)
=
\frac{27}{512\pi\,r}
\int_{[-1,1]^3}
\Bigl(\prod_{i=1}^3 (1-s_i^2)\Bigr)
K\!\left(\frac{s_1}{r},\frac{s_2}{r},\frac{s_3}{r}\right)\,ds_1ds_2ds_3.
\]

Now use the explicit formula for $D$:
\[
D\!\left(u,v;\frac{s_1}{r},\frac{s_2}{r},\frac{s_3}{r}\right)
=
\frac{1}{r}\Bigl(
s_3\sin u-s_2\sin v+s_1\sin(v-u)
\Bigr)
+ O(r^{-3}),
\]
uniformly for $u,v\in[0,2\pi]$ and $s_i\in[-1,1]$, since
\[
a_i=\sqrt{1-\frac{s_i^2}{r^2}}=1+O(r^{-2}).
\]
Therefore
\[
K\!\left(\frac{s_1}{r},\frac{s_2}{r},\frac{s_3}{r}\right)
=
\frac{1}{r}L(s_1,s_2,s_3)+O(r^{-3}),
\]
where
\[
L(s_1,s_2,s_3)
:=
\int_0^{2\pi}\int_0^{2\pi}
\left|
s_3\sin u-s_2\sin v+s_1\sin(v-u)
\right|\,du\,dv.
\]
Substituting this into the expression for $f_R(r)$ gives
\[
f_R(r)
=
\frac{A}{r^2}+O(r^{-4}),
\qquad r\to\infty,
\]
with
\[
\boxed{
A
=
\frac{27}{512\pi}
\int_{[-1,1]^3}
\Bigl(\prod_{i=1}^3 (1-s_i^2)\Bigr)
L(s_1,s_2,s_3)\,ds_1ds_2ds_3.
}
\]
Integrating the density tail, we obtain
\[
\boxed{
\mathbb{P}(R>r)=\frac{A}{r}+O(r^{-3}),
\qquad r\to\infty.
}
\]
In particular,
\[
\boxed{
\mathbb{E}[R]=\infty.
}
\]

So the heavy-tail exponent is rigorous, even though the explicit constant $A$ remains to be
evaluated.

\subsection*{5.3. Moments and Tail Dominance}

Since $\mathbb{E}[R] = \infty$, we examine fractional moments. For $0 < p < 1$,
\[
\mathbb{E}[R^p]
=
\int_0^\infty r^p f_R(r)\,dr
=
\int_0^1 r^p f_R(r)\,dr + \int_1^\infty r^p f_R(r)\,dr.
\]

The first integral is finite because $f_R(r) = O(r^2)$ as $r \to 0$. For the tail integral, using
$f_R(r) \sim A/r^2$,
\[
\int_1^\infty r^p \cdot \frac{A}{r^2}\,dr = A\int_1^\infty r^{p-2}\,dr < \infty \quad \iff \quad p < 1.
\]

Therefore $\mathbb{E}[R^p] < \infty$ for all $0 < p < 1$, with the tail contribution dominating:
\[
\int_1^\infty r^p f_R(r)\,dr \sim \frac{A}{1-p}, \quad p \nearrow 1.
\]

The variance $\mathbb{E}[R^2] = \infty$, so the distribution is heavy-tailed without finite second moment.

\section*{6. Higher-Dimensional Generalization}

The one-plane law and intersection problem generalize to $\mathbb{R}^d$ for any $d \ge 2$.

\subsection*{6.1. Planes in $\mathbb{R}^d$}

In $\mathbb{R}^d$, a $(d-2)$-dimensional affine subspace (hyperplane) is given by
\[
P = \{x \in \mathbb{R}^d : n \cdot x = h\}, \qquad n \in S^{d-1}, \quad h \in [-1,1].
\]

For three random points $x_1,x_2,x_3$ on $S^{d-1}$, the intersection
\[
C(n,h) := S^{d-1} \cap \{x : n \cdot x = h\}
\]
is a $(d-2)$-dimensional sphere of radius $\rho = \sqrt{1-h^2}$. The parametrization by angles and
the Jacobian calculation generalize: the surface-area element factors as
\[
d\sigma(x_1)\,d\sigma(x_2)\,d\sigma(x_3)
\propto
(1-h^2)^{(d-2)/2} \, dh \, d\sigma(n) \, d\Omega,
\]
where $d\Omega$ is the measure on the space of angles. After integrating out angles,
\[
\boxed{
g_d(h) = c_d(1-h^2)^{(d-2)/2}, \quad -1 \le h \le 1,
}
\]
for a normalization constant $c_d$ depending on $d$.

\subsection*{6.2. Intersection in $\mathbb{R}^d$}

For $d$ independent planes in $\mathbb{R}^d$, the intersection point is
\[
X = N^{-1}H, \quad N = \begin{pmatrix} n_1^\top \\ \vdots \\ n_d^\top \end{pmatrix}, \quad H = \begin{pmatrix} h_1 \\ \vdots \\ h_d \end{pmatrix}.
\]

By the same conditioning argument,
\[
f_X(x) = \mathbb{E}\left[ |\det(n_1,\ldots,n_d)| \prod_{i=1}^d g_d(n_i \cdot x) \right].
\]

The isotropy and radial reduction proceed as before: the radial density is
\[
f_R(r) = V_{d-1}(1) \cdot r^{d-1} f_X(re_1),
\]
where $V_{d-1}(1) = 2\pi^{d/2}/\Gamma(d/2)$ is the surface area of $S^{d-1}$.

\subsection*{6.3. Tail Behavior in Higher Dimensions}

By similar asymptotics, the tail of $f_R(r)$ scales as $r^{d-1} \cdot r^{-d} = r^{-1}$ for all $d \ge 2$.
This suggests that the heavy-tail exponent is universal in $d$: the first moment of $R$ diverges
in all dimensions. However, the constant coefficient $A_d$ and the local behavior near the origin
do depend on $d$.

\section*{7. Open Problems and Future Work}

The exact problem admits several natural extensions and open questions:

\begin{enumerate}

\item \textbf{Closed-form density:} Is there a closed-form expression for $f_R(r)$ in terms of
classical special functions (hypergeometric, Bessel, etc.)? The nested integrals suggest the
answer may be no, but a proof either way would be valuable.

\item \textbf{Evaluation of the constant $A$:} Compute the triple integral
\[
A
=
\frac{27}{512\pi}
\int_{[-1,1]^3}
(1-s_1^2)(1-s_2^2)(1-s_3^2)
L(s_1,s_2,s_3)\,ds_1ds_2ds_3
\]
where $L(s_1,s_2,s_3)$ is defined in Section 5.2. Monte Carlo simulations suggest $A \approx 2/3$.
Can this be proved exactly?

\item \textbf{Specific Probabilities:} Numerical evidence strongly suggests that the probability
of the intersection point lying outside the unit sphere is
\[
P(R > 1) = \frac{2}{\pi}.
\]
Is there a geometric reason or a simplification of the integral for $r=1$ that yields this value?
Similarly, for a single coordinate $X_i$, we observe $P(|X_i| > 1) = 1/3$.

\item \textbf{Finite-dimensional approximation of $f_R$:} Develop an asymptotic or perturbative
expansion in terms of special functions valid over the entire range $r \ge 0$.

\item \textbf{Higher moments:} Compute or bound $\mathbb{E}[R^p]$ for $0 < p < 1$ explicitly in
terms of the constant $A$.

\item \textbf{Random polytope connection:} Relate this problem to the classical random polytope
literature (e.g., the convex hull of random points). What is the combinatorial structure of
the arrangement of random planes?

\item \textbf{Non-uniform spacing:} How does the distribution of $R$ change if the three points
defining each plane are not uniformly distributed, but follow some other law on $S^2$?

\end{enumerate}

\section*{8. Conclusion}

The clean exact part of the problem is the one-plane law:
\[
g(h)=\frac{3}{4}(1-h^2)\mathbf{1}_{[-1,1]}(h).
\]
From it one gets an exact integral formula for the density of the intersection point $X$, and hence
for the distance $R=\|X\|$.

At present I do not see an elementary closed form for $f_R$. The remaining nontrivial piece is the
angular average of the absolute determinant. However, the note now gives:

\begin{itemize}
\item the exact law of a single random plane,
\item an exact integral formula for the density of the intersection point,
\item the exact local asymptotic
\[
f_R(r)=\frac{27\pi^2}{128}r^2+O(r^4),
\]
\item the rigorous tail form
\[
\mathbb{P}(R>r)=\frac{A}{r}+O(r^{-3}),
\]
with $A$ given by an explicit integral,
\item a characterization of fractional moments: $\mathbb{E}[R^p] < \infty$ for $p < 1$, with
tail dominance,
\item a natural generalization to higher dimensions $\mathbb{R}^d$, showing the offset density
scales as $(1-h^2)^{(d-2)/2}$,
\item evidence that the heavy-tail exponent $-1$ (equivalently, $\mathbb{P}(R>r) \sim A/r$) is
universal across dimensions.
\end{itemize}

\noindent
The problem connects classical random geometry with the study of random plane arrangements and
suggests further research directions both in computing the explicit constant $A$ and in exploring
variants and generalizations.

\end{document}
